\chapter{Acknowledgements}\label{ch:ack}

\begin{BlindRegion}

This template was originally developed by \gls{FCA} doctoral student \orcidlinki{Ka-Hou~Chan}{0000-0002-0183-0685}
at the request of \textit{Prof.} \orcidlinki{Chan-Tong~Lam}{0000-0002-8022-7744}
to assist \gls{PhD} students in composing their research proposals and theses in accordance with the pedagogical office’s guidelines using {\LaTeX}.
Over time, the template was refined and enhanced, ultimately resulting in a stable and reliable version.

As the institution transitioned to university. and Ka-Hou prepared to write his thesis, he redesigned the cover pages and headings, incorporating the new \gls{MPU} logo(s) to reflect the university’s identity.
Meanwhile, \textit{Prof.} \orcidlinki{Wei~Ke}{0000-0003-0952-0961} transformed the template into a {\LaTeX} class file, integrating essential commands and refining its structure to make it even more versatile and organised.

Prof. Wei Ke presented this sample thesis from the edited text of his original article titled ``Pattern Matching Based on Object Graphs''~\cite{ke2021pattern},
coauthored by Wei Ke and Ka-Hou Chan, and
published in IEEE Access, 2021 (\href{https://doi.org/10.1109/ACCESS.2021.3128575}{\sc doi: \sf\small 10.1109/ACCESS.2021.3128575}).
It is formatted with additional materials, such as source code, to demonstrate the use of the template and other commonly applicable {\LaTeX} techniques in doctoral theses.

\autoref{ch:template} lists the options and commands provided by this thesis template,
together with the history of versions and changes.
The student name ``\InsertStudentName'' and the supervisor name ``\InsertSupervisor'' are {\em fictional}. \autoref{app:biography} has more about the origins of these names.

\end{BlindRegion}
